\chapteropener{Introduction}{cyanlight}{cyanaccent}{%
% Santiago writes the opening text himself (review note, 2026-09-06). Left blank on purpose.
}{}

\clearpage
\thispagestyle{customstyle}
\section{Who this is for}

You build models, or you keep an eye on models that are already running, and you know what a confusion matrix and a residual are. That is
all we assume. What you probably don't have is a clear answer to the question that comes up in every model review: which number should
we be looking at, and what does it hide?

That is the whole point of this book. Each metric gets two pages. Not a lecture, just the things you'd want a colleague to tell you
before you pick it: what it measures, when it is the right choice, and where it will fool you.

\section{How to read a metric page}

Every metric follows the same layout, so once you've read one you know where to look on all the others.

\textbf{Page one} is the definition. A short paragraph on what the number means in practice, then the formula with every symbol
labeled, then how to read the value: the range, what counts as good, what counts as bad. After that, when to use it and when not to, and
two boxes side by side: strengths on the left, weaknesses on the right. We tried to make every bullet concrete. ``Sensitive to outliers''
tells you nothing; ``one prediction off by 100 counts as much as 99 predictions off by 1'' does.

\textbf{Page two} is the metric in action. The figure shows how the metric turns inputs into a score, and every number on it comes
from the data it draws. The caption tells you how to read it. Then a ``did you know'' box with one fact that changes how you think
about the metric, and a short note on related metrics: which neighbors to consider, and what it means when two of them disagree.

The data observability chapter is the exception. Its monitors are simple statistics, so each one gets a single page with a figure.

\section{The colors in the formulas}

Every symbol in a formula is colored by what it is, and the arrows spell it out:

\begin{itemize}[left=0pt]
    \item {\color{nmlcyan}\textbf{Cyan}} is the truth: targets, labels, actual values.
    \item {\color{nmlpurple}\textbf{Purple}} is what the model produced: predictions, scores, estimates.
    \item {\color{nmlred}\textbf{Red}} is a count or a size: how many samples, classes, bins or items.
    \item {\color{nmlgreen}\textbf{Green}} is a special function, like the indicator $\mathds{1}[\cdot]$.
    \item {\color{nmlyellow}\textbf{Yellow}} is something you choose: a threshold, a weight, a quantile level.
\end{itemize}

The same colors mean the same thing in every chapter.

\section{Picking a metric}

Three questions settle most of it.

\textbf{What does a mistake cost?} If a big miss hurts much more than a small one, you want squared errors. If a false alarm and a
missed case cost different amounts, you want precision or recall rather than accuracy, or the expected cost directly. If you can't say
what a mistake costs, the metric is deciding for you.

\textbf{How rare is the thing you care about?} Rare positives make accuracy and ROC AUC look great while the model misses them.
Tiny targets make percentage errors explode. Different base rates across groups make most fairness criteria impossible to satisfy at
the same time. Check the balance before you check the score.

\textbf{Who is going to read the number?} A metric in the target's own units, like MAE, is for the people who own the target. A
scale-free one, like $R^2$ or nDCG, is for comparing models across datasets. A proper scoring rule, like log loss, is for training and
calibration. One model can deserve all three, reported to different people.

Each chapter starts with a one-page key that walks through these questions for its metrics. And when two metrics that should agree
don't, the disagreement is usually the most useful thing on the page: RMSE far above MAE means a few large errors; ROC AUC high with
PR AUC low means the rare class is being missed.

\section{How the book is organized}

Chapters follow what the model outputs. Regression and classification first, since almost everyone meets those. Then clustering and
ranking, then computer vision, NLP and generative models, where the reference answer gets fuzzier and learned judges take over. The
last three chapters step back: estimating a metric before the labels arrive, checking whether the score is the same for every group,
and watching the tables that feed all of it. Every figure is generated by code in the book's repository, from the data it shows.
